\section{2020--2025：硅谷情绪如何从裁员转向 AI 主导}

本节先回看 2020--2025 的市场情绪变化，因为 Freda 的投资判断不是凭空出现的，而是在连续几轮利率、通胀、裁员、AI 融资和二级市场波动中形成的。

本章先建立时间背景。Freda 回顾 2020--2025 的硅谷关键词：2020 年疫情与降息，2021 年市场持续火热但中小盘先见顶，2022 年通胀和纳指大跌，2023 年 AI 在市场还没反应过来时突然成为主线，2024 年是 AI 叙事最好做的一年，2025 年则叠加关税、中美关系和继续积雪的 AI 浪潮。

这里最重要的观察，是 AI 的融资信号早于大众叙事。2022 年二级市场一片灰暗，但 Inflection AI、Anthropic、Cohere、Stability 等模型公司仍能在 A 轮融到异常大的金额。Freda 的解释很直接：Scaling Law、买卡、资本支出，这些词让 VC 的钱很早就流向了 Nvidia 和模型公司。Crossover 基金的优势在这里出现：如果一时无法判断几十家模型公司谁胜出，就可以用二级市场的 Nvidia 表达对整个 AI 需求的观点。

\begin{knowledgebox}{术语消化：Crossover、Beta 与 Alpha}
\begin{tabularx}{\textwidth}{p{0.22\textwidth}p{0.32\textwidth}X}
\toprule
术语 & 一句话解释 & 本期中的作用 \\
\midrule
Crossover 基金 & 同时投资一级私有公司和二级上市公司的基金 & 能用 Nvidia、Meta、Robinhood 等公开市场表达对前沿趋势的判断。 \\
Beta & 来自市场或行业整体上涨的收益 & 如果只判断加密周期，买 Coinbase 未必比买比特币更有 Alpha。 \\
Alpha & 超越行业/指数解释的公司自身收益 & Robinhood 被用来说明公司执行力和业务重估如何带来 Alpha。 \\
CAPEX & Capital Expenditure，资本开支 & AI 数据中心、GPU、训练和云基础设施都高度依赖 CAPEX。 \\
\bottomrule
\end{tabularx}
\end{knowledgebox}

\subsection{本章小结}

2020--2025 的硅谷不是平滑复苏，而是连续剧烈切换。AI 成为主线之前，资本已经通过大额模型公司融资和 Nvidia 需求看到了异常信号。理解这段背景，才能理解为什么本期一直在一二级市场之间切换。

